\subsubsection{运输与中继能源成本的分解}
联合方案运输能耗66.251280 kWh，中继能耗2.688136 kWh，中继占总能耗约3.90\%。表\ref{tab:rev_q3_transport_energy}按51条实际运输航段聚合，水平部分63.050347 kWh，爬升部分3.200933 kWh。爬升项在总量中所占比例较小，但仍会改变接近安全边界任务的机型和箱组可行性，因此在逐段计算中完整保留。
\begin{table}[htbp]\centering\small
\caption{联合方案运输能耗的物理分解}\label{tab:rev_q3_transport_energy}
\begin{tabular}{lrrr}\toprule
机型 & 水平能耗/kWh & 爬升能耗/kWh & 总能耗/kWh\\\midrule
A & 24.888150 & 1.516529 & 26.404679\\
B & 14.308489 & 0.814536 & 15.123025\\
C & 23.853709 & 0.869867 & 24.723576\\
\bottomrule\end{tabular}
\end{table}


四项中继任务的转场合计0.866261 kWh，建链和实际服务合计1.821876 kWh，后者约占中继能耗的67.77\%。表\ref{tab:rev_q3_relay_energy}显示，各任务服务10.780—34.583 min，占自身准备至返航作业时间的44.78\%—66.96\%。中继成本不仅取决于飞到哪里，还取决于在那里等待和服务多久。
\begin{table}[htbp]\centering\small
\caption{中继任务的转场、建链与悬停服务成本}\label{tab:rev_q3_relay_energy}
\begin{tabular}{lrrrrr}\toprule
中继任务 & 转场/kWh & 建链及服务/kWh & 总能耗/kWh & 服务/min & 服务占用比/\%\\\midrule
Q3-R-01 & 0.225372 & 0.373222 & 0.598594 & 19.858 & 52.37\\
Q3-R-02 & 0.210978 & 0.643196 & 0.854174 & 34.583 & 66.96\\
Q3-R-03 & 0.142681 & 0.206793 & 0.349474 & 10.780 & 44.78\\
Q3-R-04 & 0.287229 & 0.598665 & 0.885894 & 32.154 & 60.38\\
\bottomrule\end{tabular}
\end{table}


能源充足和服务及时是两种条件。中继返航SOC较高，提供了能源上的余量；但如果它未在关键切换前完成建链，运输任务仍可能等待。因而改善一个方案应先识别起作用的约束：位置可能主要改善遮挡和交接，开始时刻可能主要减少提前悬停，组件配置则主要影响跨架次恢复。这些作用都能在分解表与实际服务窗口中对应。

\subsubsection{通信提供方的实际保障工作量}
从正式保存的181个通信阶段、228条实际提供方区间重新聚合，运输任务累计需要34312.495780 s保障。固定网关承担20734.854799 s，占60.43\%；中继承担13577.640980 s，占39.57\%。18个运输架次在部分阶段使用中继，5个全程直连，见表\ref{tab:rev_q3_provider}。
\begin{table}[htbp]\centering\small
\caption{实际通信提供方的累计保障工作量}\label{tab:rev_q3_provider}
\begin{tabular}{lrrrr}\toprule
提供方 & 区间数 & 运输任务数 & 累计保障/min & 比例/\%\\\midrule
G01 & 112 & 23 & 345.581 & 60.43\\
Q3-R-01 & 41 & 7 & 92.738 & 16.22\\
Q3-R-02 & 28 & 7 & 48.093 & 8.41\\
Q3-R-03 & 15 & 3 & 24.757 & 4.33\\
Q3-R-04 & 32 & 9 & 60.706 & 10.62\\
\bottomrule\end{tabular}
\end{table}


累计保障时长按运输任务分别求和。如果两架运输机在同一秒使用同一中继，该秒贡献两份保障工作量。因此Q3-R-01约19.858 min的在线窗口，可对应约92.738 min的多任务累计保障。该指标说明其承担的任务服务量，不是系统失联时长、通信吞吐量或中继自己的飞行时长。当前模型没有额外的并发用户容量约束，保障关系按给定双跳链路条件核算。

实际提供方序列在合并相邻同提供方记录后共有47次切换，包括直连—中继和中继—中继交接。表\ref{tab:rev_q3_comm_tasks}列出中继保障需求较大的任务。相对于静态“覆盖多少服务区”，这些记录更直接回答：中继为哪些运输机服务、在什么时段服务、哪些切换会影响任务开始。
\begin{table}[htbp]\centering\small
\caption{中继保障需求较大的运输任务}\label{tab:rev_q3_comm_tasks}
\begin{tabular}{lrrr}\toprule
运输任务 & 直连/min & 中继/min & 提供方切换数\\\midrule
Q3-T-09 & 13.587 & 24.505 & 3\\
Q3-T-16 & 14.928 & 21.465 & 3\\
Q3-T-17 & 16.300 & 20.589 & 3\\
Q3-T-12 & 11.503 & 16.227 & 3\\
Q3-T-20 & 12.676 & 13.521 & 2\\
Q3-T-10 & 12.263 & 13.283 & 3\\
Q3-T-21 & 12.676 & 13.021 & 2\\
Q3-T-15 & 21.829 & 12.739 & 4\\
\bottomrule\end{tabular}
\end{table}


中继以较小的能源占比支撑较大比例的运输通信任务时间，体现了运输与通信资源的分工价值。该统计不替代优化证明，而为联合方案增加运行机理说明：空间部署提供可用链路，服务窗口保证时间匹配，实体和组件复用保证这些服务能够实际安排。
